% 第18章题目解答
% 按 chapters/ch18.tex 中出现顺序编排；普通例题不重复作答。

\chapter*{第18章\quad 多元函数的极限与连续习题解答}
\addcontentsline{toc}{chapter}{第18章\quad 多元函数的极限与连续习题解答}
\subsection*{18.1.4 思考题}

\paragraph{1. 多元函数极限的 Heine 归结原理}
\begin{xsolution}
先证必要性。设
\[
  \lim_{\vect{x}\to\vect{a}}f(\vect{x})=A,
\]
并取任意满足 $\vect{x}_n\ne\vect{a}$、$\vect{x}_n\to\vect{a}$ 的点列。给定 $\varepsilon>0$，由极限定义，存在 $\delta>0$，使得
\[
  0<\abs{\vect{x}-\vect{a}}<\delta
  \Longrightarrow
  \abs{f(\vect{x})-A}<\varepsilon.
\]
由于 $\vect{x}_n\to\vect{a}$，存在 $N$，当 $n\ge N$ 时
\[
  0<\abs{\vect{x}_n-\vect{a}}<\delta,
\]
于是
\[
  \abs{f(\vect{x}_n)-A}<\varepsilon.
\]
故 $f(\vect{x}_n)\to A$。

再证充分性。反设 $\lim_{\vect{x}\to\vect{a}}f(\vect{x})\ne A$。则存在 $\varepsilon_0>0$，使得对任意 $\delta>0$，总能找到定义域中的点 $\vect{x}$ 满足
\[
  0<\abs{\vect{x}-\vect{a}}<\delta,
  \qquad
  \abs{f(\vect{x})-A}\ge\varepsilon_0.
\]
令 $\delta=1/n$，依次选取 $\vect{x}_n$，使
\[
  0<\abs{\vect{x}_n-\vect{a}}<\frac1n,
  \qquad
  \abs{f(\vect{x}_n)-A}\ge\varepsilon_0.
\]
于是 $\vect{x}_n\to\vect{a}$，但 $f(\vect{x}_n)$ 不趋于 $A$，与题设矛盾。因此
\[
  \lim_{\vect{x}\to\vect{a}}f(\vect{x})=A.
\]
\end{xsolution}

\paragraph{2. 多元函数极限的 Cauchy 收敛准则}
\begin{xsolution}
必要性。设
\[
  \lim_{\vect{x}\to\vect{a}}f(\vect{x})=A.
\]
给定 $\varepsilon>0$，取 $\delta>0$，使
\[
  0<\abs{\vect{x}-\vect{a}}<\delta
  \Longrightarrow
  \abs{f(\vect{x})-A}<\frac\varepsilon2.
\]
于是对任意 $\vect{x}',\vect{x}''\in O_\delta(\vect{a})\setminus\{\vect{a}\}$，
\[
  \abs{f(\vect{x}')-f(\vect{x}'')}
  \le
  \abs{f(\vect{x}')-A}+\abs{f(\vect{x}'')-A}
  <\varepsilon.
\]

充分性。设题中 Cauchy 条件成立。由于 $\vect{a}$ 是定义域的聚点，可取定义域中的点列 $\vect{x}_n\ne\vect{a}$，使 $\vect{x}_n\to\vect{a}$。给定 $\varepsilon>0$，取 Cauchy 条件中的 $\delta>0$。当 $m,n$ 足够大时，
\[
  \vect{x}_m,\vect{x}_n\in O_\delta(\vect{a})\setminus\{\vect{a}\},
\]
所以
\[
  \abs{f(\vect{x}_m)-f(\vect{x}_n)}<\varepsilon.
\]
故 $\{f(\vect{x}_n)\}$ 是实数 Cauchy 列，从而收敛，记
\[
  f(\vect{x}_n)\to A.
\]

任取另一列 $\vect{y}_n\ne\vect{a}$、$\vect{y}_n\to\vect{a}$，把两列交错组成
\[
  \vect{z}_{2n}=\vect{x}_n,
  \qquad
  \vect{z}_{2n-1}=\vect{y}_n.
\]
仍有 $\vect{z}_n\to\vect{a}$，由同一 Cauchy 条件知 $\{f(\vect{z}_n)\}$ 是 Cauchy 列，故收敛。其偶数项子列已经趋于 $A$，所以整列极限只能是 $A$，于是奇数项子列也满足
\[
  f(\vect{y}_n)\to A.
\]
由上一题的 Heine 归结原理，得到
\[
  \lim_{\vect{x}\to\vect{a}}f(\vect{x})=A.
\]
\end{xsolution}

\subsection*{18.1.6 练习题}

\paragraph{1. 求下列极限}

\textbf{(1)}
\begin{xsolution}
令
\[
  s(t)=
  \begin{cases}
    \dfrac{\sin t}{t}, & t\ne0,\\[4pt]
    1, & t=0.
  \end{cases}
\]
由基本极限知 $s$ 在 $t=0$ 连续。对题中函数的定义域 $x\ne0$，无论 $y$ 是否为 $0$，均有
\[
  \frac{\sin xy}{x}=y s(xy).
\]
当 $(x,y)\to(0,a)$ 时，$xy\to0$，故 $s(xy)\to1$，同时 $y\to a$。因此
\[
  \boxed{\displaystyle
  \lim_{(x,y)\to(0,a)}\frac{\sin xy}{x}=a}.
\]
\end{xsolution}

\textbf{(2)}
\begin{xsolution}
令 $r=\sqrt{x^2+y^2}$。由 Cauchy--Schwarz 不等式，
\[
  \abs{x+y}\le\sqrt2\,r.
\]
于是当 $0<r<1$ 时，
\[
  \abs{(x+y)\ln(x^2+y^2)}
  \le
  \sqrt2\,r\abs{\ln r^2}
  =2\sqrt2\,r\abs{\ln r}.
\]
而 $r\abs{\ln r}\to0$（$r\to0^+$），故
\[
  \boxed{\displaystyle
  \lim_{(x,y)\to(0,0)}(x+y)\ln(x^2+y^2)=0}.
\]
\end{xsolution}

\textbf{(3)}
\begin{xsolution}
在 $t=0$ 的邻近定义
\[
  \ell(t)=
  \begin{cases}
    \dfrac{\ln(1+t)}{t}, & t\ne0,\\[4pt]
    1, & t=0.
  \end{cases}
\]
由基本极限知 $\ell$ 在 $0$ 连续。于是对 $x\ne0$ 且 $1+xy>0$，包括 $y=0$ 的情形，都有
\[
  f(x,y)=y\ell(xy).
\]
当 $(x,y)\to(0,0)$ 时，$xy\to0$，故 $\ell(xy)$ 在充分小邻域内有界：存在 $C>0$ 使
\[
  \abs{\ell(xy)}\le C.
\]
因此 $x\ne0$ 时 $\abs{f(x,y)}\le C\abs{y}$；而 $x=0$ 时 $\abs{f(0,y)}=\abs{y}$。取 $C_1\ge\max\{C,1\}$，则原点附近统一有
\[
  \abs{f(x,y)}\le C_1\abs{y}\to0.
\]
所以
\[
  \boxed{\displaystyle
  \lim_{(x,y)\to(0,0)}f(x,y)=0}.
\]
\end{xsolution}

\textbf{(4)}
\begin{xsolution}
在底数为正的充分远邻域内取对数，记
\[
  L(x,y)
  =\frac{x^2}{x+y}\ln\left(1+\frac1x\right)
  =\frac{x}{x+y}\,x\ln\left(1+\frac1x\right).
\]
当 $(x,y)\to(\infty,0)$ 时，
\[
  \frac{x}{x+y}=\frac1{1+y/x}\to1,
\]
且
\[
  x\ln\left(1+\frac1x\right)\to1
  \qquad(\abs{x}\to\infty).
\]
故 $L(x,y)\to1$，从而
\[
  \boxed{\displaystyle
  \lim_{(x,y)\to(\infty,0)}
  \left(1+\frac1x\right)^{\frac{x^2}{x+y}}
  =\ee}.
\]
\end{xsolution}

\textbf{(5)}
\begin{xsolution}
有
\[
  x^2-xy+y^2
  =\frac12(x^2+y^2)+\frac12(x-y)^2
  \ge\frac12(x^2+y^2),
\]
并且
\[
  \abs{x+y}\le\sqrt2\sqrt{x^2+y^2}.
\]
因此
\[
  \abs*{\frac{x+y}{x^2-xy+y^2}}
  \le
  \frac{2\sqrt2}{\sqrt{x^2+y^2}}\to0
\]
当 $(x,y)\to(\infty,\infty)$ 时成立。故
\[
  \boxed{\displaystyle
  \lim_{(x,y)\to(\infty,\infty)}
  \frac{x+y}{x^2-xy+y^2}=0}.
\]
\end{xsolution}

\paragraph{2. 证明：$\displaystyle\lim_{(x,y)\to(0,0)}\frac{x^2y^2}{x^3+y^3}$ 不存在}
\begin{xsolution}
沿曲线
\[
  y=-x+cx^2,
  \qquad c>0,
\]
趋于原点。此时
\[
  x^3+y^3=(x+y)(x^2-xy+y^2),
\]
而
\[
  x+y=cx^2,
\]
\[
  x^2-xy+y^2=x^2(3-3cx+c^2x^2),
\]
\[
  x^2y^2=x^4(1-cx)^2.
\]
故在该曲线上
\[
  \frac{x^2y^2}{x^3+y^3}
  =\frac{(1-cx)^2}{c(3-3cx+c^2x^2)}
  \longrightarrow\frac1{3c}.
\]
取 $c=1$ 与 $c=2$，分别得到 $1/3$ 与 $1/6$。两条趋近曲线所得极限不同，所以
\[
  \boxed{\displaystyle
  \lim_{(x,y)\to(0,0)}\frac{x^2y^2}{x^3+y^3}
  \text{不存在}}.
\]
\end{xsolution}

\paragraph{3. 讨论下列函数在点 $(0,0)$ 的重极限与累次极限}

\textbf{(1)}
\begin{xsolution}
设
\[
  f(x,y)=\frac{x^2y^2}{x^2y^2+(x-y)^2}.
\]
沿 $y=x$ 且 $x\ne0$，有 $f(x,x)=1$；沿 $y=0$，有 $f(x,0)=0$。故重极限不存在。

对固定 $y\ne0$，当 $x\to0$ 时 $f(x,y)\to0$，所以
\[
  \lim_{y\to0}\lim_{x\to0}f(x,y)=0.
\]
同理，对固定 $x\ne0$，当 $y\to0$ 时 $f(x,y)\to0$，故
\[
  \lim_{x\to0}\lim_{y\to0}f(x,y)=0.
\]
因此两个累次极限均存在且等于 $0$，但重极限不存在。
\end{xsolution}

\textbf{(2)}
\begin{xsolution}
设
\[
  f(x,y)=(x+y)\sin\frac1x\sin\frac1y.
\]
在定义域内
\[
  \abs{f(x,y)}\le\abs{x}+\abs{y}\to0,
\]
故
\[
  \boxed{\displaystyle
  \lim_{(x,y)\to(0,0)}f(x,y)=0}.
\]

若固定 $y\ne0$ 且 $\sin(1/y)\ne0$，则
\[
  f(x,y)
  =x\sin\frac1x\sin\frac1y
  +y\sin\frac1y\sin\frac1x.
\]
第一项趋于 $0$，第二项因 $\sin(1/x)$ 振荡而无极限，故 $\lim_{x\to0}f(x,y)$ 对任意小邻域中的许多 $y$ 不存在。因此
\[
  \lim_{y\to0}\lim_{x\to0}f(x,y)
\]
不存在。交换 $x,y$ 后同理可知
\[
  \lim_{x\to0}\lim_{y\to0}f(x,y)
\]
也不存在。
\end{xsolution}

\textbf{(3)}
\begin{xsolution}
当 $x\ne0$ 时
\[
  f(x,y)=\frac{xy}{x^2+y^2}+y\sin\frac1x,
\]
而 $x=0$ 时 $f(0,y)=0$。

沿 $y=0$ 有 $f(x,0)=0$；沿 $y=x$ 有
\[
  f(x,x)=\frac12+x\sin\frac1x\to\frac12.
\]
故重极限不存在。

先令 $y\to0$。对固定 $x\ne0$，
\[
  \frac{xy}{x^2+y^2}\to0,
  \qquad
  y\sin\frac1x\to0,
\]
而 $x=0$ 时函数恒为 $0$，所以
\[
  \lim_{y\to0}f(x,y)=0.
\]
于是
\[
  \boxed{\displaystyle
  \lim_{x\to0}\lim_{y\to0}f(x,y)=0}.
\]

反之，固定 $y\ne0$。当 $x\to0$ 时第一项趋于 $0$，但 $y\sin(1/x)$ 振荡，故 $\lim_{x\to0}f(x,y)$ 不存在。因此
\[
  \boxed{\displaystyle
  \lim_{y\to0}\lim_{x\to0}f(x,y)
  \text{不存在}}.
\]
\end{xsolution}

\paragraph{4. 设 $g(x,y)=\dfrac{f(x)-f(y)}{x-y}$，求 $\displaystyle\lim_{(x,y)\to(t,t)}g(x,y)$}
\begin{xsolution}
对 $x\ne y$，由微积分基本定理，
\[
  f(x)-f(y)=\int_y^x f'(s)\,\dd s.
\]
令 $s=y+u(x-y)$，得
\[
  g(x,y)
  =\int_0^1f'\bigl(y+u(x-y)\bigr)\,\dd u.
\]
当 $(x,y)\to(t,t)$ 时，$y+u(x-y)\to t$ 关于 $u\in[0,1]$ 一致。由 $f'$ 在 $t$ 连续，对任意 $\varepsilon>0$，当 $(x,y)$ 足够接近 $(t,t)$ 时，
\[
  \abs{f'\bigl(y+u(x-y)\bigr)-f'(t)}<\varepsilon,
  \qquad0\le u\le1.
\]
于是
\[
  \abs{g(x,y)-f'(t)}
  \le
  \int_0^1\abs{f'\bigl(y+u(x-y)\bigr)-f'(t)}\,\dd u
  <\varepsilon.
\]
因此
\[
  \boxed{\displaystyle
  \lim_{(x,y)\to(t,t)}g(x,y)=f'(t)}.
\]

\textbf{另解。} 对 $x\ne y$，由一元函数的中值定理，存在介于 $x$ 与 $y$ 之间的点 $\xi$，使
\[
  g(x,y)=\frac{f(x)-f(y)}{x-y}=f'(\xi).
\]
当 $(x,y)\to(t,t)$ 时，$\xi\to t$。由 $f'$ 连续，立即得到 $g(x,y)\to f'(t)$。
\end{xsolution}

\paragraph{5. 叙述并证明二元函数极限的惟一性、局部有界性与局部保号性}
\begin{xsolution}
设 $f(x,y)$ 在点 $(x_0,y_0)$ 的某个去心邻域内有定义。

\textbf{惟一性。} 若
\[
  \lim_{(x,y)\to(x_0,y_0)}f(x,y)=A,
  \qquad
  \lim_{(x,y)\to(x_0,y_0)}f(x,y)=B,
\]
则 $A=B$。

反设 $A\ne B$，取 $\varepsilon=\abs{A-B}/3$。由两个极限的定义，在同一个足够小的去心邻域内同时有
\[
  \abs{f(x,y)-A}<\varepsilon,
  \qquad
  \abs{f(x,y)-B}<\varepsilon.
\]
于是
\[
  \abs{A-B}
  \le\abs{A-f(x,y)}+\abs{f(x,y)-B}
  <2\varepsilon
  =\frac23\abs{A-B},
\]
矛盾。

\textbf{局部有界性。} 若
\[
  \lim_{(x,y)\to(x_0,y_0)}f(x,y)=A\in\RR,
\]
则存在 $\delta>0$ 与 $M>0$，使
\[
  0<\sqrt{(x-x_0)^2+(y-y_0)^2}<\delta
  \Longrightarrow
  \abs{f(x,y)}\le M.
\]
取 $\varepsilon=1$。在某个去心邻域内
\[
  \abs{f(x,y)-A}<1,
\]
故
\[
  \abs{f(x,y)}\le\abs{A}+1.
\]
取 $M=\abs{A}+1$ 即可。

\textbf{局部保号性。} 若
\[
  \lim_{(x,y)\to(x_0,y_0)}f(x,y)=A\ne0,
\]
则在某个去心邻域内 $f(x,y)$ 与 $A$ 同号，且
\[
  \abs{f(x,y)}>\frac{\abs{A}}{2}.
\]
事实上，取 $\varepsilon=\abs{A}/2$。在充分小的去心邻域内
\[
  \abs{f(x,y)-A}<\frac{\abs{A}}{2}.
\]
若 $A>0$，则 $f(x,y)>A/2>0$；若 $A<0$，则 $f(x,y)<A/2<0$。同时
\[
  \abs{f(x,y)}
  \ge\abs{A}-\abs{f(x,y)-A}
  >\frac{\abs{A}}{2}.
\]
\end{xsolution}

\paragraph{6. 由 $\varphi(y)$ 与 $\psi(x)$ 的控制证明重极限}
% SOURCE-ERR: 按通常的二元函数极限定义，原题条件不足；沿直线 $y=y_0$ 或 $x=x_0$ 时，一元极限条件并不控制 $\varphi(y_0)$、$\psi(x_0)$。
\begin{xsolution}
按题面现有条件，结论不成立。取 $x_0=y_0=0$、$A=0$，定义
\[
  \varphi(y)=
  \begin{cases}
    0, & y\ne0,\\
    1, & y=0,
  \end{cases}
  \qquad
  \psi(x)=
  \begin{cases}
    0, & x\ne0,\\
    1, & x=0,
  \end{cases}
\]
以及
\[
  f(x,y)=\varphi(y).
\]
则
\[
  \lim_{y\to0}\varphi(y)=0,
  \qquad
  \lim_{x\to0}\psi(x)=0,
\]
而且处处有
\[
  \abs{f(x,y)-\varphi(y)}=0\le\psi(x).
\]
但沿 $y=0$、$x\to0$ 有 $f(x,0)=1$，所以
\[
  \lim_{(x,y)\to(0,0)}f(x,y)
\]
不存在，更不等于 $A=0$。

若补充
\[
  \varphi(y_0)=A,
  \qquad
  \psi(x_0)=0,
\]
则原结论成立。事实上，给定 $\varepsilon>0$，由两个一元极限分别取 $\delta_1,\delta_2>0$，使在相应邻域内
\[
  \abs{\varphi(y)-A}<\frac\varepsilon2,
  \qquad
  0\le\psi(x)<\frac\varepsilon2,
\]
其中 $y=y_0$ 与 $x=x_0$ 的情形由补充条件保证。于是当 $(x,y)$ 充分接近 $(x_0,y_0)$ 时，
\[
  \abs{f(x,y)-A}
  \le\abs{f(x,y)-\varphi(y)}+\abs{\varphi(y)-A}
  \le\psi(x)+\abs{\varphi(y)-A}
  <\varepsilon.
\]
因此在上述补充条件下才有
\[
  \lim_{(x,y)\to(x_0,y_0)}f(x,y)=A.
\]
\end{xsolution}

\paragraph{7. 命题 18.1.2 中补充定义域条件后的结论}
\begin{xsolution}
沿用命题 18.1.2 的记号。由该命题，存在同一个 $A$，使
\[
  \lim_{x\to x_0}g(x)=A,
  \qquad
  \lim_{y\to y_0}h(y)=A.
\]
条件 (2) 表明
\[
  \lim_{x\to x_0}f(x,y)=h(y)
\]
在 $0<\abs{y-y_0}<\eta$ 上关于 $y$ 一致。因此给定 $\varepsilon>0$，存在 $\delta_1>0$，使
\[
  0<\abs{x-x_0}<\delta_1,
  \quad
  0<\abs{y-y_0}<\eta
\]
时恒有
\[
  \abs{f(x,y)-h(y)}<\frac\varepsilon2.
\]
又由 $h(y)\to A$，存在 $\delta_2\in(0,\eta)$，使
\[
  0<\abs{y-y_0}<\delta_2
  \Longrightarrow
  \abs{h(y)-A}<\frac\varepsilon2.
\]
于是当
\[
  0<\abs{x-x_0}<\delta_1,
  \qquad
  0<\abs{y-y_0}<\delta_2
\]
时，
\[
  \abs{f(x,y)-A}
  \le\abs{f(x,y)-h(y)}+\abs{h(y)-A}
  <\varepsilon.
\]
题设的补充条件说明函数在点 $(x_0,y_0)$ 邻近除两条直线 $x=x_0$、$y=y_0$ 外均有定义；上面的估计覆盖了定义域中所有充分接近 $(x_0,y_0)$ 的点，因此
\[
  \lim_{(x,y)\to(x_0,y_0)}f(x,y)=A.
\]
两个累次极限分别为
\[
  \lim_{x\to x_0}\lim_{y\to y_0}f(x,y)
  =\lim_{x\to x_0}g(x)=A,
\]
\[
  \lim_{y\to y_0}\lim_{x\to x_0}f(x,y)
  =\lim_{y\to y_0}h(y)=A.
\]
故重极限与两个累次极限都存在且相等。
\end{xsolution}

\subsection*{18.2.5 练习题}

\paragraph{1. 证明 $F_n(x)=f(x,\varphi_n(x))$ 一致收敛}
\begin{xsolution}
设
\[
  \varphi_n\to\varphi
\]
在 $[a,b]$ 上一致。由 $c\le\varphi_n(x)\le d$，令 $n\to\infty$ 得
\[
  c\le\varphi(x)\le d.
\]
由于 $f$ 在紧集 $[a,b]\times[c,d]$ 上连续，故一致连续。给定 $\varepsilon>0$，存在 $\delta>0$，使得同属该矩形的两点 $(x,y_1),(x,y_2)$ 只要满足
\[
  \abs{y_1-y_2}<\delta,
\]
就有
\[
  \abs{f(x,y_1)-f(x,y_2)}<\varepsilon.
\]
由 $\varphi_n\to\varphi$ 一致，存在 $N$，当 $n\ge N$ 时
\[
  \sup_{x\in[a,b]}\abs{\varphi_n(x)-\varphi(x)}<\delta.
\]
于是
\[
  \sup_{x\in[a,b]}
  \abs{f(x,\varphi_n(x))-f(x,\varphi(x))}<\varepsilon.
\]
所以 $F_n$ 在 $[a,b]$ 上一致收敛到
\[
  F(x)=f(x,\varphi(x)).
\]
\end{xsolution}

\paragraph{2. 证明 $f(x,y)=\sqrt{x^2+y^2}$ 在 $\RR^2$ 上一致连续}
\begin{xsolution}
把两点记为 $\vect{p}=(x_1,y_1)$、$\vect{q}=(x_2,y_2)$。由三角不等式的反向形式，
\[
  \abs*{\sqrt{x_1^2+y_1^2}-\sqrt{x_2^2+y_2^2}}
  =\bigl|\abs{\vect{p}}-\abs{\vect{q}}\bigr|
  \le\abs{\vect{p}-\vect{q}}.
\]
因此该函数满足 Lipschitz 条件，常数为 $1$。给定 $\varepsilon>0$，取 $\delta=\varepsilon$ 即得一致连续性。
\end{xsolution}

\paragraph{3. 命题 18.2.1 中由 (1) 直接证明 (4)}
\begin{xsolution}
设 $f$ 在 $\RR^n$ 上连续，$E\subset\RR^n$。任取
\[
  \vect{x}\in\overline E.
\]
存在点列 $\vect{x}_k\in E$，使 $\vect{x}_k\to\vect{x}$。由连续性，
\[
  f(\vect{x}_k)\to f(\vect{x}).
\]
而 $f(\vect{x}_k)\in f(E)$，故
\[
  f(\vect{x})\in\overline{f(E)}.
\]
由于 $\vect{x}\in\overline E$ 任意，得到
\[
  \boxed{f(\overline E)\subset\overline{f(E)}}.
\]
\end{xsolution}

\paragraph{4. 如果 $f$ 把 $\RR$ 中任意开集映为开集，问 $f$ 是否连续}
\begin{xsolution}
不一定。

把 $\RR$ 按有理数加法子群 $\QQ$ 的陪集分解。每个陪集
\[
  x+\QQ
\]
都是 $\RR$ 中的稠密集，这些陪集彼此不交，且陪集族与 $\RR$ 等势。取一个双射
\[
  \beta\colon\RR/\QQ\to\RR,
\]
定义
\[
  f(x)=\beta(x+\QQ).
\]
任意非空开集 $U\subset\RR$ 与每个陪集都相交，所以
\[
  f(U)=\RR,
\]
是开集；空集的像仍为空集。因此 $f$ 把任意开集映为开集。

但 $f$ 处处不连续。事实上，固定 $x_0$，任取一个实数 $c$ 使
\[
  \abs{c-f(x_0)}>1.
\]
令 $C=\beta^{-1}(c)$ 为相应陪集。由于 $C$ 稠密，在 $x_0$ 的任意邻域中都能取 $x\in C$，于是
\[
  \abs{f(x)-f(x_0)}=\abs{c-f(x_0)}>1.
\]
故 $f$ 在 $x_0$ 不连续。由 $x_0$ 任意，$f$ 处处不连续。
\end{xsolution}

\paragraph{5. 连续函数在无穷远有有限极限时有界且一致连续}
\begin{xsolution}
设
\[
  \lim_{\abs{x}+\abs{y}\to\infty}f(x,y)=A.
\]

先证有界。取 $R>0$，使
\[
  \abs{x}+\abs{y}>R
  \Longrightarrow
  \abs{f(x,y)-A}<1.
\]
故在该区域
\[
  \abs{f(x,y)}\le\abs{A}+1.
\]
而紧集
\[
  K_R=\{(x,y)\mid\abs{x}+\abs{y}\le R\}
\]
上 $f$ 连续，故有界。合并两部分，$f$ 在 $\RR^2$ 上有界。

再证一致连续。给定 $\varepsilon>0$，取 $R>0$，使
\[
  \abs{x}+\abs{y}>R
  \Longrightarrow
  \abs{f(x,y)-A}<\frac\varepsilon2.
\]
在紧集
\[
  K=\{(x,y)\mid\abs{x}+\abs{y}\le R+2\}
\]
上，$f$ 一致连续。因此存在 $\delta_0>0$，使 $\vect{p},\vect{q}\in K$ 且
\[
  \abs{\vect{p}-\vect{q}}<\delta_0
\]
时，有
\[
  \abs{f(\vect{p})-f(\vect{q})}<\varepsilon.
\]
取
\[
  \delta=\min\left\{\delta_0,\frac1{\sqrt2}\right\}.
\]
若 $\abs{\vect{p}-\vect{q}}<\delta$，且两点均在 $K$ 中，结论已由紧集上的一致连续性得到。

若至少一点不在 $K$，不妨
\[
  \abs{p_1}+\abs{p_2}>R+2.
\]
则
\begin{align*}
  \abs{q_1}+\abs{q_2}
  &\ge
  \abs{p_1}+\abs{p_2}
  -\abs{p_1-q_1}-\abs{p_2-q_2}\\
  &>R+2-\sqrt2\delta
  \ge R+1>R.
\end{align*}
故两点都在无穷远区域，因而
\[
  \abs{f(\vect{p})-f(\vect{q})}
  \le\abs{f(\vect{p})-A}+\abs{f(\vect{q})-A}
  <\varepsilon.
\]
所以 $f$ 在 $\RR^2$ 上一致连续。
\end{xsolution}

\paragraph{6. 若 $D\subset\RR^2$ 是有界闭域，则 $f(D)$ 是有界闭区间}
\begin{xsolution}
有界闭域 $D$ 是紧集，并且作为“域”是连通的。连续函数 $f$ 把紧集映为紧集，所以 $f(D)\subset\RR$ 是有界闭集；又由连续映射保持连通性，$f(D)$ 是连通集。$\RR$ 中的连通集必为区间。因此 $f(D)$ 是一个有界闭区间。
\end{xsolution}

\paragraph{7. 利用命题 18.2.1 和例题 18.2.7 重新证明上一章第一组参考题 5（$\RR^n$ 的正规性）}
% SOURCE-ERR: chapters/ch18.tex 此处作“上一章的第二组参考题 5”；扫描原页明确作“上一章的第一组参考题 5”。
\begin{xsolution}
设 $S_1,S_2$ 是 $\RR^n$ 中两个不相交的闭集。令
\[
  d_i(\vect{x})=d(\vect{x},S_i),
  \qquad i=1,2.
\]
由例题 18.2.7，$d_1,d_2$ 都在 $\RR^n$ 上连续。因为 $S_i$ 闭，
\[
  d_i(\vect{x})=0
  \Longleftrightarrow
  \vect{x}\in S_i.
\]
又因 $S_1\cap S_2=\varnothing$，对任意 $\vect{x}\in\RR^n$，$d_1(\vect{x})$ 与 $d_2(\vect{x})$ 不会同时为 $0$。因此
\[
  h(\vect{x})
  =\frac{d_1(\vect{x})}{d_1(\vect{x})+d_2(\vect{x})}
\]
是 $\RR^n$ 上的连续函数。

令
\[
  O_1=h^{-1}\bigl((-\infty,1/3)\bigr),
  \qquad
  O_2=h^{-1}\bigl((2/3,+\infty)\bigr).
\]
由命题 18.2.1，开集的原象为开集，故 $O_1,O_2$ 都是开集。由于区间 $(-\infty,1/3)$ 与 $(2/3,+\infty)$ 不相交，故
\[
  O_1\cap O_2=\varnothing.
\]
若 $\vect{x}\in S_1$，则 $h(\vect{x})=0$，所以 $S_1\subset O_1$；若 $\vect{x}\in S_2$，则 $h(\vect{x})=1$，所以 $S_2\subset O_2$。因此存在互不相交的开集 $O_1,O_2$ 分别包含 $S_1,S_2$，即 $\RR^n$ 是正规空间。
\end{xsolution}

\paragraph{8. Dini 判别法在紧集上的推广}
\begin{xsolution}
由
\[
  f_1(\vect{x})\ge f_2(\vect{x})\ge\cdots,
  \qquad
  f_k(\vect{x})\to0,
\]
可知对每个 $\vect{x}\in D$ 及每个 $k$，均有 $f_k(\vect{x})\ge0$。

给定 $\varepsilon>0$，令
\[
  G_k=\{\vect{x}\in D\mid f_k(\vect{x})<\varepsilon\}.
\]
由于 $f_k$ 连续，$G_k$ 是 $D$ 中的相对开集；又因 $f_{k+1}\le f_k$，有
\[
  G_1\subset G_2\subset\cdots.
\]
对每个 $\vect{x}\in D$，由 $f_k(\vect{x})\to0$，存在 $k$ 使 $f_k(\vect{x})<\varepsilon$，故
\[
  D=\bigcup_{k=1}^\infty G_k.
\]
由 $D$ 紧，可从该开覆盖中取有限子覆盖
\[
  D=G_{k_1}\cup\cdots\cup G_{k_m}.
\]
取 $N=\max\{k_1,\ldots,k_m\}$。由于 $G_k$ 递增，得到
\[
  D=G_N.
\]
于是对一切 $\vect{x}\in D$，
\[
  0\le f_N(\vect{x})<\varepsilon.
\]
当 $k\ge N$ 时，$0\le f_k\le f_N$，故
\[
  \sup_{\vect{x}\in D}\abs{f_k(\vect{x})}<\varepsilon.
\]
这正说明 $f_k\to0$ 在 $D$ 上一致。
\end{xsolution}

\paragraph{9. 例题 18.2.8 中振幅的另一种表达}
\begin{xsolution}
定义
\[
  \widetilde\omega_f(\vect{a},\delta)
  =\sup_{\substack{\vect{x}\in E\\
  \abs{\vect{a}-\vect{x}}\le\delta}}
  \abs{f(\vect{a})-f(\vect{x})}.
\]
当 $\delta$ 增大时，取上确界的集合增大，所以 $\widetilde\omega_f(\vect{a},\delta)$ 关于 $\delta$ 单调增加。若 $f$ 在 $E$ 上有界，则该量有界，因此
\[
  \widetilde\omega_f(\vect{a})
  =\lim_{\delta\to0^+}\widetilde\omega_f(\vect{a},\delta)
\]
存在。

因为 $f(\vect{a})$ 本身也是邻域中的一个函数值，
\[
  \widetilde\omega_f(\vect{a},\delta)
  \le\omega_f(\vect{a},\delta).
\]
另一方面，对邻域中的任意 $\vect{x},\vect{y}$，
\[
  \abs{f(\vect{x})-f(\vect{y})}
  \le
  \abs{f(\vect{x})-f(\vect{a})}
  +\abs{f(\vect{y})-f(\vect{a})}
  \le2\widetilde\omega_f(\vect{a},\delta).
\]
取上确界得到
\[
  \omega_f(\vect{a},\delta)
  \le2\widetilde\omega_f(\vect{a},\delta).
\]
令 $\delta\to0^+$，有
\[
  \widetilde\omega_f(\vect{a})
  \le\omega_f(\vect{a})
  \le2\widetilde\omega_f(\vect{a}).
\]
因此
\[
  \omega_f(\vect{a})=0
  \Longleftrightarrow
  \widetilde\omega_f(\vect{a})=0.
\]
结合例题 18.2.8，$f$ 在 $\vect{a}$ 连续当且仅当
\[
  \lim_{\delta\to0^+}
  \sup_{\substack{\vect{x}\in E\\
  \abs{\vect{a}-\vect{x}}\le\delta}}
  \abs{f(\vect{a})-f(\vect{x})}=0.
\]
这里“可改用”指它给出同一个连续性判据；一般而言，两种振幅的数值不必完全相等。
\end{xsolution}

\paragraph{10. 证明连续映射把紧集映为紧集}
\begin{xsolution}
设 $K\subset\RR^n$ 为紧集，$f\colon K\to\RR^m$ 连续。任取 $f(K)$ 中的点列
\[
  \vect{y}_j=f(\vect{x}_j),
  \qquad \vect{x}_j\in K.
\]
由 $K$ 紧，$\{\vect{x}_j\}$ 有收敛子列
\[
  \vect{x}_{j_k}\to\vect{x}_0\in K.
\]
由连续性，
\[
  \vect{y}_{j_k}
  =f(\vect{x}_{j_k})
  \to f(\vect{x}_0)\in f(K).
\]
因此 $f(K)$ 中任意点列都有收敛于 $f(K)$ 中一点的子列，即 $f(K)$ 紧。
\end{xsolution}

\subsection*{18.3.2 第一组参考题}

\paragraph{1. 单调参数下的边界收敛关于 $x\in[a,b]$ 一致}
\begin{xsolution}
由于 $f(x,y)$ 关于 $y$ 单调增加，且
\[
  \lim_{y\to d^-}f(x,y)=f(x,d),
\]
所以对 $y<d$ 有
\[
  f(x,y)\le f(x,d).
\]
给定 $\varepsilon>0$。对每个 $x_0\in[a,b]$，由点态收敛可取 $y_{x_0}<d$，使
\[
  0\le f(x_0,d)-f(x_0,y_{x_0})<\frac\varepsilon3.
\]
函数 $x\mapsto f(x,d)$ 与 $x\mapsto f(x,y_{x_0})$ 都在 $[a,b]$ 上连续，因此存在 $x_0$ 的邻域 $U_{x_0}$，使对 $x\in U_{x_0}\cap[a,b]$，
\[
  \abs{f(x,d)-f(x_0,d)}<\frac\varepsilon3,
\]
\[
  \abs{f(x,y_{x_0})-f(x_0,y_{x_0})}<\frac\varepsilon3.
\]
于是
\[
  0\le f(x,d)-f(x,y_{x_0})<\varepsilon.
\]

这些 $U_{x_0}$ 覆盖紧集 $[a,b]$，可取有限子覆盖
\[
  U_{x_1},\ldots,U_{x_m}.
\]
令
\[
  y_*=\max_{1\le i\le m}y_{x_i},
  \qquad
  \delta=d-y_*>0.
\]
若 $0<d-y<\delta$，则 $y>y_*\ge y_{x_i}$。对任意 $x\in[a,b]$，选 $i$ 使 $x\in U_{x_i}$。由关于 $y$ 的单调性，
\[
  0\le f(x,d)-f(x,y)
  \le f(x,d)-f(x,y_{x_i})<\varepsilon.
\]
所以
\[
  \sup_{x\in[a,b]}\abs{f(x,y)-f(x,d)}<\varepsilon,
\]
即收敛关于 $x\in[a,b]$ 一致。
\end{xsolution}

\paragraph{2. 证明 $\varphi(x)$ 在 $[a,b]$ 上连续}
\begin{xsolution}
把约束区域参数化。令
\[
  \xi=a+s(x-a),
  \qquad
  y=a+t(\xi-a)=a+ts(x-a),
  \qquad 0\le s,t\le1.
\]
则所有满足 $a\le y\le\xi\le x$ 的点都可如此表示，而且
\[
  \varphi(x)
  =\max_{0\le s,t\le1}
  f\bigl(a+s(x-a),a+ts(x-a)\bigr).
\]
设
\[
  F(x,s,t)=f\bigl(a+s(x-a),a+ts(x-a)\bigr).
\]
$F$ 在紧集 $[a,b]\times[0,1]\times[0,1]$ 上连续，故一致连续。

给定 $x_0\in[a,b]$ 与 $\varepsilon>0$。当 $x$ 足够接近 $x_0$ 时，对所有 $s,t\in[0,1]$ 都有
\[
  \abs{F(x,s,t)-F(x_0,s,t)}<\varepsilon.
\]
于是
\[
  F(x,s,t)\le F(x_0,s,t)+\varepsilon
  \le\varphi(x_0)+\varepsilon.
\]
对 $s,t$ 取最大值得
\[
  \varphi(x)\le\varphi(x_0)+\varepsilon.
\]
交换 $x,x_0$ 同理得到
\[
  \varphi(x_0)\le\varphi(x)+\varepsilon.
\]
故
\[
  \abs{\varphi(x)-\varphi(x_0)}\le\varepsilon.
\]
所以 $\varphi$ 在 $[a,b]$ 上连续。
\end{xsolution}

\paragraph{3. 紧集上的函数连续当且仅当其图像是紧集}
\begin{xsolution}
记函数图像为
\[
  \Gamma_f
  =\{(\vect{x},y)\in\RR^{n+1}\mid
  \vect{x}\in D,\ y=f(\vect{x})\}.
\]

若 $f$ 在紧集 $D$ 上连续，则映射
\[
  \Phi\colon D\to\RR^{n+1},
  \qquad
  \Phi(\vect{x})=(\vect{x},f(\vect{x}))
\]
连续。连续映射把紧集映为紧集，故
\[
  \Gamma_f=\Phi(D)
\]
是紧集。

反之，设 $\Gamma_f$ 紧。任取 $\vect{x}_k\to\vect{x}\in D$。若 $f(\vect{x}_k)\not\to f(\vect{x})$，则存在 $\varepsilon_0>0$ 与子列 $\vect{x}_{k_j}$，使
\[
  \abs{f(\vect{x}_{k_j})-f(\vect{x})}\ge\varepsilon_0.
\]
由于
\[
  (\vect{x}_{k_j},f(\vect{x}_{k_j}))\in\Gamma_f
\]
且 $\Gamma_f$ 紧，可再取子列使
\[
  (\vect{x}_{k_{j_\ell}},f(\vect{x}_{k_{j_\ell}}))
  \to(\vect{u},v)\in\Gamma_f.
\]
第一分量已经趋于 $\vect{x}$，故 $\vect{u}=\vect{x}$。极限点属于图像，必须有
\[
  v=f(\vect{x}).
\]
于是
\[
  f(\vect{x}_{k_{j_\ell}})\to f(\vect{x}),
\]
与上面的 $\varepsilon_0$ 不等式矛盾。因此对任意 $\vect{x}_k\to\vect{x}$ 都有 $f(\vect{x}_k)\to f(\vect{x})$，由 Heine 判别法，$f$ 在 $D$ 上连续。
\end{xsolution}

\paragraph{4. 连续函数组线性相关与 Gram 行列式}
\begin{xsolution}
令 Gram 矩阵
\[
  \mat{G}=(G_{ij})_{n\times n},
  \qquad
  G_{ij}=\int_0^1f_i(x)f_j(x)\,\dd x.
\]

若 $f_1,\ldots,f_n$ 线性相关，则存在非零向量
\[
  \vect{c}=(c_1,\ldots,c_n)^{\mathsf T}
\]
使
\[
  \sum_{j=1}^nc_jf_j(x)=0,
  \qquad0\le x\le1.
\]
于是对每个 $i$，
\[
  (\mat{G}\vect{c})_i
  =\sum_{j=1}^nG_{ij}c_j
  =\int_0^1f_i(x)\sum_{j=1}^nc_jf_j(x)\,\dd x
  =0.
\]
故 $\mat{G}\vect{c}=0$ 且 $\vect{c}\ne0$，从而 $\det\mat{G}=0$。

反之，若 $\det\mat{G}=0$，则存在 $\vect{c}\ne0$ 使 $\mat{G}\vect{c}=0$。于是
\[
  0=\vect{c}^{\mathsf T}\mat{G}\vect{c}
  =\int_0^1
  \left(\sum_{j=1}^nc_jf_j(x)\right)^2\,\dd x.
\]
被积函数连续且非负，其积分为 $0$，故它在 $[0,1]$ 上恒为 $0$。因此
\[
  \sum_{j=1}^nc_jf_j(x)=0,
  \qquad0\le x\le1,
\]
且 $\vect{c}\ne0$，所以 $f_1,\ldots,f_n$ 线性相关。
\end{xsolution}

\paragraph{5. 有限点集的最小半径覆盖圆盘}
\begin{xsolution}
对圆心 $\vect{c}\in\RR^2$ 定义
\[
  R(\vect{c})
  =\max_{1\le i\le k}\abs{\vect{c}-p_i}.
\]
这是有限个连续函数的最大值，故连续。设
\[
  M_0=\max_{1\le i\le k}\abs{p_i}.
\]
由三角不等式，
\[
  R(\vect{c})\ge\abs{\vect{c}}-M_0,
\]
所以
\[
  R(\vect{c})\to\infty
  \qquad(\abs{\vect{c}}\to\infty).
\]

取固定圆心 $p_1$，令 $R_0=R(p_1)$。当
\[
  \abs{\vect{c}}>R_0+M_0
\]
时，$R(\vect{c})>R_0$。因此全空间上的最小值一定可在紧球
\[
  \{\vect{c}\mid\abs{\vect{c}}\le R_0+M_0\}
\]
中寻找。由连续函数在紧集上达到最小值，存在 $\vect{c}_*$ 使
\[
  R(\vect{c}_*)=\min_{\vect{c}\in\RR^2}R(\vect{c}).
\]
以 $\vect{c}_*$ 为圆心、$R(\vect{c}_*)$ 为半径的闭圆盘覆盖所有 $p_i$，且半径最小。

在 $\RR^n$ 中完全相同，只需把圆盘改为闭球，仍令
\[
  R(\vect{c})=\max_i\abs{\vect{c}-p_i}.
\]
上述连续性与趋于无穷的论证不变。
\end{xsolution}

\paragraph{6. 广义 Rayleigh 商}

\textbf{(1)}
\begin{xsolution}
因为 $\mat{B}$ 正定，
\[
  \vect{x}^{\mathsf T}\mat{B}\vect{x}>0,
  \qquad\vect{x}\ne\vect{0}.
\]
考虑椭球面
\[
  S_{\mat{B}}
  =\{\vect{x}\in\RR^n\mid
  \vect{x}^{\mathsf T}\mat{B}\vect{x}=1\}.
\]
二次型连续，故 $S_{\mat{B}}$ 闭。正定二次型在单位球面上有正的最小值 $m>0$，所以
\[
  \vect{x}^{\mathsf T}\mat{B}\vect{x}
  \ge m\abs{\vect{x}}^2.
\]
从而 $\vect{x}\in S_{\mat{B}}$ 时
\[
  \abs{\vect{x}}\le\frac1{\sqrt m}.
\]
故 $S_{\mat{B}}$ 紧。

在 $S_{\mat{B}}$ 上
\[
  G(\vect{x})=\vect{x}^{\mathsf T}\mat{A}\vect{x}
\]
连续，因此达到最大值。任意 $\vect{x}\ne0$ 可归一化为
\[
  \vect{y}
  =\frac{\vect{x}}
  {\sqrt{\vect{x}^{\mathsf T}\mat{B}\vect{x}}}
  \in S_{\mat{B}},
\]
且齐次性给出
\[
  G(\vect{x})=G(\vect{y}).
\]
所以 $G$ 在整个 $E=\RR^n\setminus\{\vect{0}\}$ 上也能取到最大值。
\end{xsolution}

\textbf{(2)}
\begin{xsolution}
设 $\vect{x}_*$ 是最大值点，并按齐次性归一化使
\[
  \vect{x}_*^{\mathsf T}\mat{B}\vect{x}_*=1.
\]
于是 $\vect{x}_*$ 是约束
\[
  \vect{x}^{\mathsf T}\mat{B}\vect{x}=1
\]
下二次型 $\vect{x}^{\mathsf T}\mat{A}\vect{x}$ 的极大点。由 Lagrange 乘子法，存在 $\lambda\in\RR$，使
\[
  2\mat{A}\vect{x}_*=2\lambda\mat{B}\vect{x}_*,
\]
即
\[
  \mat{A}\vect{x}_*=\lambda\mat{B}\vect{x}_*.
\]
由于 $\mat{B}$ 可逆，
\[
  \mat{B}^{-1}\mat{A}\vect{x}_*=\lambda\vect{x}_*.
\]
所以最大值点 $\vect{x}_*$ 是矩阵
\[
  \boxed{\mat{B}^{-1}\mat{A}}
\]
的特征向量。左乘 $\vect{x}_*^{\mathsf T}$ 并利用归一化条件可得
\[
  \lambda
  =\vect{x}_*^{\mathsf T}\mat{A}\vect{x}_*
  =G(\vect{x}_*),
\]
即对应特征值就是广义 Rayleigh 商的最大值。
\end{xsolution}

\subsection*{18.3.2 第二组参考题}

\paragraph{1. 若 $T^{n_0}$ 是压缩映射，则 $T$ 有惟一不动点}
\begin{xsolution}
令
\[
  S=T^{n_0}.
\]
题设说明 $S$ 是 $\RR^n$ 上的压缩映射。由压缩映射原理，存在惟一的 $\vect{x}_*$ 使
\[
  S\vect{x}_*=\vect{x}_*.
\]
又因为 $S$ 与 $T$ 交换，
\[
  ST=TS,
\]
所以
\[
  S(T\vect{x}_*)=T(S\vect{x}_*)=T\vect{x}_*.
\]
故 $T\vect{x}_*$ 也是 $S$ 的不动点。由 $S$ 的不动点惟一性，
\[
  T\vect{x}_*=\vect{x}_*.
\]
所以 $\vect{x}_*$ 是 $T$ 的不动点。

若 $\vect{y}$ 也是 $T$ 的不动点，则
\[
  S\vect{y}=T^{n_0}\vect{y}=\vect{y},
\]
故 $\vect{y}$ 也是 $S$ 的不动点，于是 $\vect{y}=\vect{x}_*$。因此 $T$ 有惟一不动点。
\end{xsolution}

\paragraph{2. 紧集上的严格缩短映射}
\begin{xsolution}
由
\[
  \abs{f(\vect{x})-f(\vect{y})}<\abs{\vect{x}-\vect{y}},
  \qquad\vect{x}\ne\vect{y},
\]
可知 $f$ 连续。定义
\[
  h(\vect{x})=\abs{f(\vect{x})-\vect{x}}.
\]
$\Omega$ 有界闭，故紧；$h$ 连续，所以存在 $\vect{x}_*\in\Omega$ 使
\[
  h(\vect{x}_*)=\min_{\vect{x}\in\Omega}h(\vect{x}).
\]
若 $h(\vect{x}_*)>0$，则 $f(\vect{x}_*)\ne\vect{x}_*$，且 $f(\vect{x}_*)\in\Omega$。于是
\[
  h(f(\vect{x}_*))
  =\abs{f(f(\vect{x}_*))-f(\vect{x}_*)}
  <\abs{f(\vect{x}_*)-\vect{x}_*}
  =h(\vect{x}_*),
\]
与最小性矛盾。因此
\[
  f(\vect{x}_*)=\vect{x}_*.
\]
若有两个不同的不动点 $\vect{x}_*,\vect{y}_*$，则
\[
  \abs{\vect{x}_*-\vect{y}_*}
  =\abs{f(\vect{x}_*)-f(\vect{y}_*)}
  <\abs{\vect{x}_*-\vect{y}_*},
\]
矛盾，故不动点惟一。

有界闭集的条件不能分别减弱为一般有界集或一般无界闭集。

对一般有界集，取
\[
  \Omega=(0,1),
  \qquad
  f(x)=\frac x2.
\]
则 $f(\Omega)\subset\Omega$ 且
\[
  \abs{f(x)-f(y)}=\frac12\abs{x-y}<\abs{x-y},
\]
但惟一候选不动点 $0\notin\Omega$，故没有不动点。

对无界闭集，取
\[
  \Omega=[0,+\infty),
  \qquad
  f(x)=x+\ee^{-x}.
\]
对 $x\ge0$ 有 $f(x)=x+\ee^{-x}>x\ge0$，所以 $f(\Omega)\subset\Omega$，且等式 $f(x)=x$ 无解，即没有不动点。若 $x>y\ge0$，由中值定理存在 $\xi\in(y,x)$ 使
\[
  f(x)-f(y)=f'(\xi)(x-y),
\]
而
\[
  f'(\xi)=1-\ee^{-\xi}\in[0,1),
\]
故
\[
  \abs{f(x)-f(y)}<\abs{x-y}.
\]
因此无界闭集也不足以保证结论。
\end{xsolution}

\paragraph{3. 证明连续映射将连通集映为连通集}
\begin{xsolution}
设 $D$ 连通，$f\colon D\to Y$ 连续。反设 $f(D)$ 不连通，则存在 $f(D)$ 中两个不相交的非空相对开集 $A,B$，使
\[
  f(D)=A\cup B.
\]
于是
\[
  D=f^{-1}(A)\cup f^{-1}(B).
\]
由于 $f$ 连续，$f^{-1}(A)$ 与 $f^{-1}(B)$ 都是 $D$ 中的相对开集；它们非空、不相交，且并为 $D$。这与 $D$ 连通矛盾。因此 $f(D)$ 连通。
\end{xsolution}

\paragraph{4. 上、下半连续的极限刻画}

\textbf{(1) 上半连续}
\begin{xsolution}
先设 $f$ 在 $\vect{x}_0$ 上半连续。取 $\varepsilon=1$，存在 $\delta_0>0$，使
\[
  \vect{x}\in\Omega,
  \quad
  \abs{\vect{x}-\vect{x}_0}<\delta_0
  \Longrightarrow
  f(\vect{x})<f(\vect{x}_0)+1.
\]
故 $f$ 在 $\vect{x}_0$ 邻近有上界。

给定任意 $\varepsilon>0$。由上半连续性，在充分小邻域内
\[
  f(\vect{x})<f(\vect{x}_0)+\varepsilon.
\]
所以
\[
  \limsup_{\Omega\ni\vect{x}\to\vect{x}_0}f(\vect{x})
  \le f(\vect{x}_0)+\varepsilon.
\]
令 $\varepsilon\downarrow0$，得
\[
  \limsup_{\Omega\ni\vect{x}\to\vect{x}_0}f(\vect{x})
  \le f(\vect{x}_0).
\]

反之，设 $f$ 在 $\vect{x}_0$ 邻近有上界，且
\[
  L:=\limsup_{\Omega\ni\vect{x}\to\vect{x}_0}f(\vect{x})
  \le f(\vect{x}_0).
\]
给定 $\varepsilon>0$，由上极限定义，存在 $\delta>0$，使
\[
  \sup_{\substack{\vect{x}\in\Omega\\
  0<\abs{\vect{x}-\vect{x}_0}<\delta}}
  f(\vect{x})
  <L+\varepsilon
  \le f(\vect{x}_0)+\varepsilon.
\]
把 $\vect{x}=\vect{x}_0$ 包括进来后，题中上半连续不等式仍成立，所以 $f$ 在 $\vect{x}_0$ 上半连续。
\end{xsolution}

\textbf{(2) 下半连续}
\begin{xsolution}
若 $f$ 在 $\vect{x}_0$ 下半连续，取 $\varepsilon=1$ 得到邻近有下界。对任意 $\varepsilon>0$，在充分小邻域内
\[
  f(\vect{x})>f(\vect{x}_0)-\varepsilon.
\]
故
\[
  \liminf_{\Omega\ni\vect{x}\to\vect{x}_0}f(\vect{x})
  \ge f(\vect{x}_0)-\varepsilon.
\]
令 $\varepsilon\downarrow0$，得到
\[
  \liminf_{\Omega\ni\vect{x}\to\vect{x}_0}f(\vect{x})
  \ge f(\vect{x}_0).
\]

反之，若 $f$ 在邻近有下界且
\[
  L:=\liminf_{\Omega\ni\vect{x}\to\vect{x}_0}f(\vect{x})
  \ge f(\vect{x}_0),
\]
则给定 $\varepsilon>0$，由下极限定义，存在 $\delta>0$，使
\[
  \inf_{\substack{\vect{x}\in\Omega\\
  0<\abs{\vect{x}-\vect{x}_0}<\delta}}
  f(\vect{x})
  >L-\varepsilon
  \ge f(\vect{x}_0)-\varepsilon.
\]
因此在充分小邻域内
\[
  f(\vect{x})>f(\vect{x}_0)-\varepsilon,
\]
即 $f$ 在 $\vect{x}_0$ 下半连续。
\end{xsolution}

\paragraph{5. 紧集上的半连续函数达到单侧极值}
\begin{xsolution}
只证上半连续情形；下半连续情形对 $-f$ 应用同一论证即可。

先证有上界。对每个 $\vect{x}\in\Omega$，由上半连续性取邻域 $U_{\vect{x}}$，使
\[
  \vect{y}\in U_{\vect{x}}\cap\Omega
  \Longrightarrow
  f(\vect{y})<f(\vect{x})+1.
\]
这些邻域覆盖紧集 $\Omega$，可取有限子覆盖
\[
  U_{\vect{x}_1},\ldots,U_{\vect{x}_m}.
\]
于是对所有 $\vect{y}\in\Omega$，
\[
  f(\vect{y})
  \le\max_{1\le i\le m}\bigl(f(\vect{x}_i)+1\bigr),
\]
故 $f$ 有上界。

令
\[
  M=\sup_{\vect{x}\in\Omega}f(\vect{x}).
\]
取点列 $\vect{x}_n\in\Omega$ 使
\[
  f(\vect{x}_n)>M-\frac1n.
\]
由紧性，可取子列 $\vect{x}_{n_k}\to\vect{x}_*\in\Omega$。上半连续的极限刻画给出
\[
  \limsup_{k\to\infty}f(\vect{x}_{n_k})
  \le f(\vect{x}_*).
\]
左端等于 $M$，所以
\[
  M\le f(\vect{x}_*)\le M.
\]
故 $f(\vect{x}_*)=M$，即达到最大值。

若 $f$ 下半连续，则 $-f$ 上半连续，故 $-f$ 达到最大值，等价于 $f$ 有下界并达到最小值。
\end{xsolution}

\paragraph{6. 证明 $f(\vect{x})=\sup_k\{f_k(\vect{x})\}$ 在 $\Omega$ 上达到最小值}
\begin{xsolution}
题设保证该上确界对每个 $\vect{x}$ 都是有限实数。对任意 $a\in\RR$，
\[
  \{\vect{x}\in\Omega\mid f(\vect{x})>a\}
  =\bigcup_{k=1}^\infty
  \{\vect{x}\in\Omega\mid f_k(\vect{x})>a\}.
\]
每个 $f_k$ 连续，故右端每一集合都开，因此其并集开。于是 $f$ 是下半连续函数。

由上一题，紧集上的下半连续函数有下界并达到最小值。因此存在 $\vect{x}_*\in\Omega$，使
\[
  f(\vect{x}_*)
  =\min_{\vect{x}\in\Omega}f(\vect{x}).
\]
\end{xsolution}

\paragraph{7. Peano 曲线}

\textbf{(1)}
\begin{xsolution}
固定 $m$。按构造，$f_m(t)$ 所在的第 $m$ 级小正三角形边长为 $2^{-m}$；之后每一步变换都保持该参数点的像留在这个第 $m$ 级小三角形中。因此对任意 $n>m$，$f_n(t)$ 与 $f_m(t)$ 位于同一个边长为 $2^{-m}$ 的小正三角形中。该三角形的直径就是边长，所以
\[
  \abs{f_n(t)-f_m(t)}\le\frac1{2^m},
  \qquad0\le t\le1.
\]
于是
\[
  \sup_{t\in[0,1]}\abs{f_n(t)-f_m(t)}
  \le\frac1{2^m},
  \qquad n>m.
\]
故 $\{f_n\}$ 是一致 Cauchy 列。$\RR^2$ 完备，所以存在映射 $f$ 使 $f_n\to f$ 在 $[0,1]$ 上一致。

每个 $f_n$ 都是有限条线段拼成的连续折线映射，连续函数列的一致极限连续，因此
\[
  \boxed{f\colon[0,1]\to\Delta\text{ 连续}}.
\]
令 $n\to\infty$ 于前面的估计，还得到
\[
  \sup_{t\in[0,1]}\abs{f(t)-f_m(t)}
  \le\frac1{2^m}.
\]
\end{xsolution}

\textbf{(2)}
\begin{xsolution}
由每个 $f_m([0,1])\subset\Delta$ 且闭三角形 $\Delta$ 为闭集，一致极限逐点仍落在 $\Delta$ 中，因此 $f([0,1])\subset\Delta$。下面证明反向包含。

任取 $p\in\Delta$。第 $m$ 次等分把 $\Delta$ 分成边长为 $2^{-m}$ 的小正三角形。取其中一个包含 $p$ 的小三角形，记为 $\Delta_m(p)$。按构造，折线 $f_m([0,1])$ 在每个第 $m$ 级小三角形中都有一段，因此可取 $t_m\in[0,1]$，使
\[
  f_m(t_m)\in\Delta_m(p).
\]
同一小三角形内两点距离不超过其直径 $2^{-m}$，故
\[
  \abs{f_m(t_m)-p}\le\frac1{2^m}.
\]
由 (1) 的一致估计，
\[
  \abs{f(t_m)-f_m(t_m)}\le\frac1{2^m}.
\]
所以
\[
  \abs{f(t_m)-p}\le\frac2{2^m}\to0.
\]

区间 $[0,1]$ 紧，故 $\{t_m\}$ 有收敛子列
\[
  t_{m_j}\to t_*\in[0,1].
\]
由 $f$ 连续，
\[
  f(t_{m_j})\to f(t_*).
\]
另一方面，前面的估计给出 $f(t_{m_j})\to p$，故
\[
  f(t_*)=p.
\]
因此任意 $p\in\Delta$ 都属于 $f([0,1])$，从而
\[
  \boxed{f([0,1])=\Delta}.
\]
正三角形边长为 $1$，故面积为 $\sqrt3/4$。
\end{xsolution}

\paragraph{8. Tietze（蒂策）扩张定理}

\textbf{(1)}
% SOURCE-ERR: 原题的 g_k 公式缺少因子 M，且距离项顺序应为 d(x,B_k)-d(x,A_k)；否则不能得到 g|_D=f。
\begin{xsolution}
题面所列递推式需要两处校正。若按题面原式，
\[
  \abs{g_k(\vect{x})}\le\frac{2^{k-1}}{3^k},
\]
故若原式处处有定义，则
\[
  \sum_{k=1}^\infty\abs{g_k(\vect{x})}
  \le\sum_{k=1}^\infty\frac{2^{k-1}}{3^k}=1.
\]
因此当 $M>1$ 且 $f$ 在某点取值大于 $1$ 时，原式不可能满足 $\sum g_k=f$。此外，若 $\vect{x}\in A_k$，原题所写公式给出的校正方向与减小正残差所需方向相反。

与 $A_k,B_k$ 的定义相容的递推式应为
\[
  \widetilde g_k(\vect{x})
  =a_kM\,
  \frac{d(\vect{x},B_k)-d(\vect{x},A_k)}
  {d(\vect{x},A_k)+d(\vect{x},B_k)},
  \qquad
  a_k=\frac{2^{k-1}}{3^k},
\]
当 $A_k,B_k$ 都非空时使用此式；若 $A_k=\varnothing,B_k\ne\varnothing$，取 $\widetilde g_k\equiv-a_kM$；若 $B_k=\varnothing,A_k\ne\varnothing$，取 $\widetilde g_k\equiv a_kM$；若二者皆空，取 $\widetilde g_k\equiv0$。以下证明这个修正式确实给出题目所要的结论。

令
\[
  \varphi_1=f,
  \qquad
  \varphi_k
  =f-\sum_{j=1}^{k-1}\widetilde g_j,
  \qquad k\ge2,
\]
在 $D$ 上定义。归纳证明
\[
  \abs{\varphi_k(\vect{x})}
  \le\left(\frac23\right)^{k-1}M
  =3a_kM,
  \qquad \vect{x}\in D.
\]
$k=1$ 时就是题设 $\abs{f}\le M$。

假设该估计对 $k$ 成立。由归纳假设，$\widetilde g_1,\ldots,\widetilde g_{k-1}$ 连续，故 $\varphi_k=f-\sum_{j=1}^{k-1}\widetilde g_j$ 在 $D$ 上连续。于是 $A_k,B_k$ 是 $D$ 中的闭集；又因 $D$ 本身闭，它们也是 $\RR^n$ 中的闭集。两集合由正、负阈值的定义互不相交。若二者非空，则距离函数连续且分母处处为正，所以 $\widetilde g_k$ 连续，并且
\[
  \abs{\widetilde g_k(\vect{x})}\le a_kM.
\]
在 $D$ 上分三种情况。

若 $\varphi_k(\vect{x})\ge a_kM$，则 $\vect{x}\in A_k$，因而
\[
  \widetilde g_k(\vect{x})=a_kM.
\]
结合 $\varphi_k\le3a_kM$ 得
\[
  0\le\varphi_{k+1}(\vect{x})
  =\varphi_k(\vect{x})-a_kM
  \le2a_kM.
\]
若 $\varphi_k(\vect{x})\le-a_kM$，则
\[
  \widetilde g_k(\vect{x})=-a_kM,
\]
从而
\[
  -2a_kM\le\varphi_{k+1}(\vect{x})\le0.
\]
若
\[
  \abs{\varphi_k(\vect{x})}<a_kM,
\]
则仅用 $\abs{\widetilde g_k}\le a_kM$ 即得
\[
  \abs{\varphi_{k+1}(\vect{x})}\le2a_kM.
\]
空集情形采用上面的常值约定时，同样得到这个估计。因此
\[
  \abs{\varphi_{k+1}(\vect{x})}
  \le2a_kM
  =\left(\frac23\right)^kM.
\]
归纳完成。

又因为
\[
  \sup_{\vect{x}\in\RR^n}
  \abs{\widetilde g_k(\vect{x})}
  \le a_kM,
\]
且
\[
  \sum_{k=1}^\infty a_kM=M<\infty,
\]
由 Weierstrass 判别法，级数
\[
  \sum_{k=1}^\infty\widetilde g_k(\vect{x})
\]
在 $\RR^n$ 上一致收敛。每一项连续，所以其和
\[
  g(\vect{x})=\sum_{k=1}^\infty\widetilde g_k(\vect{x})
\]
连续。

对 $\vect{x}\in D$，由残差估计
\[
  \abs*{f(\vect{x})-
  \sum_{j=1}^{k-1}\widetilde g_j(\vect{x})}
  =\abs{\varphi_k(\vect{x})}
  \le\left(\frac23\right)^{k-1}M\to0.
\]
故
\[
  g(\vect{x})=f(\vect{x}),
  \qquad\vect{x}\in D.
\]
这证明了修正递推式下的有界 Tietze 扩张结论。
\end{xsolution}

\textbf{(2)}
\begin{xsolution}
使用 (1) 中已经得到的有界 Tietze 扩张结论。设 $f\colon D\to\RR$ 连续，但不要求有界。定义
\[
  h(\vect{x})
  =\frac{f(\vect{x})}{1+\abs{f(\vect{x})}},
  \qquad\vect{x}\in D.
\]
则 $h$ 连续且 $\abs{h(\vect{x})}<1$。由有界情形，存在连续函数
\[
  H\colon\RR^n\to[-1,1]
\]
使 $H|_D=h$。

令
\[
  Z=\{\vect{x}\in\RR^n\mid\abs{H(\vect{x})}=1\}.
\]
$Z$ 是闭集，且因 $\abs{h}<1$，有 $D\cap Z=\varnothing$。若 $Z=\varnothing$，令 $\lambda\equiv1$；若 $Z\ne\varnothing$，定义
\[
  \lambda(\vect{x})
  =\frac{d(\vect{x},Z)}
  {d(\vect{x},Z)+d(\vect{x},D)}.
\]
两个距离函数连续。由于 $D,Z$ 闭且不交，分母处处为正，所以 $\lambda$ 连续，并且
\[
  0\le\lambda\le1,
  \qquad
  \lambda|_D=1,
  \qquad
  \lambda|_Z=0.
\]
令
\[
  K(\vect{x})=\lambda(\vect{x})H(\vect{x}).
\]
则 $K$ 连续，$K|_D=h$，并且对每个 $\vect{x}\in\RR^n$ 都有 $\abs{K(\vect{x})}<1$：若 $\vect{x}\in Z$，则 $K(\vect{x})=0$；若 $\vect{x}\notin Z$，则 $\abs{H}<1$ 且 $0\le\lambda\le1$。

最后定义
\[
  G(\vect{x})
  =\frac{K(\vect{x})}{1-\abs{K(\vect{x})}}.
\]
因为 $\abs{K}<1$，$G$ 在全空间连续。对 $\vect{x}\in D$，$K=h$，而函数
\[
  u\longmapsto\frac{u}{1-\abs{u}}
\]
是 $t\mapsto t/(1+\abs{t})$ 在 $(-1,1)$ 上的反函数，所以
\[
  G(\vect{x})=f(\vect{x}).
\]
因此无界连续函数同样可以从闭集 $D$ 连续扩张到 $\RR^n$。
\end{xsolution}

\paragraph{9. 开集上的连续函数能否扩张到全空间？一致连续函数呢？}
\begin{xsolution}
一般连续函数不能。取
\[
  U=(0,1)\subset\RR,
  \qquad
  f(x)=\frac1x.
\]
$f$ 在 $U$ 上连续。若存在连续扩张 $F\colon\RR\to\RR$，则由 $x\to0^+$ 应有
\[
  F(x)=\frac1x\to F(0)\in\RR,
\]
但 $1/x\to+\infty$，矛盾。因此开集上的连续函数未必能连续扩张到全空间。

若 $f$ 在开集 $U\subset\RR$ 上一致连续，则可以。先把它惟一连续扩张到闭包 $\overline U$。任取 $x\in\overline U$，取 $x_n\in U$、$x_n\to x$。由于 $f$ 一致连续，$\{f(x_n)\}$ 是 Cauchy 列，故收敛；一致连续性还保证其极限与所选逼近列无关。定义
\[
  \widetilde f(x)=\lim_{n\to\infty}f(x_n).
\]
同一个 $\varepsilon$--$\delta$ 条件传到极限后表明 $\widetilde f$ 在 $\overline U$ 上连续，并且 $\widetilde f|_U=f$。

现在 $\RR\setminus\overline U$ 是若干互不相交开区间的并。在每个有限空隙 $(a,b)$ 上，用连接 $\widetilde f(a)$ 与 $\widetilde f(b)$ 的线性函数定义扩张；在无界空隙 $(-\infty,a)$ 或 $(a,+\infty)$ 上取相邻端点值的常值函数。每个空隙内部的连续性直接成立，在空隙端点处也与 $\widetilde f$ 的端点值吻合。若有若干空隙的端点向 $x_0\in\overline U$ 聚集，则 $\widetilde f$ 在 $\overline U$ 上仍一致连续；足够靠近 $x_0$ 的空隙两端函数值都靠近 $\widetilde f(x_0)$，而线性插值位于两端函数值之间，故这些空隙内部的值也靠近 $\widetilde f(x_0)$。因此所得函数在全 $\RR$ 上连续，并在 $U$ 上与 $f$ 一致。

\textbf{另解（一致连续情形）。} 仍先由一致连续性把 $f$ 惟一连续扩张到闭集 $\overline U$。由于 $\overline U\subset\RR$ 是闭集，再应用上一题的 Tietze 扩张定理，即可把 $\widetilde f$ 连续扩张到整个 $\RR$。这直接得到所需的全空间连续扩张。

所以在 $\RR$ 中：一般连续函数的答案是否定的；一致连续函数的答案是肯定的。
\end{xsolution}

\paragraph{10. 有理点集上的一致连续函数可惟一扩张到全空间}
\begin{xsolution}
设有理点集为
\[
  D=\QQ^n\subset\RR^n,
\]
且 $f\colon D\to\RR$ 一致连续。任取 $\vect{x}\in\RR^n$，选有理点列
\[
  \vect{q}_m\in\QQ^n,
  \qquad
  \vect{q}_m\to\vect{x}.
\]
因为 $\{\vect{q}_m\}$ 是 Cauchy 列，而 $f$ 一致连续，所以 $\{f(\vect{q}_m)\}$ 也是 Cauchy 列，在 $\RR$ 中收敛。定义
\[
  F(\vect{x})=\lim_{m\to\infty}f(\vect{q}_m).
\]

先证定义与逼近列无关。若另有
\[
  \vect{r}_m\in\QQ^n,
  \qquad
  \vect{r}_m\to\vect{x},
\]
则
\[
  \abs{\vect{q}_m-\vect{r}_m}\to0.
\]
给定 $\varepsilon>0$，由一致连续性存在 $\delta>0$，使
\[
  \abs{\vect{u}-\vect{v}}<\delta
  \Longrightarrow
  \abs{f(\vect{u})-f(\vect{v})}<\varepsilon
\]
对所有有理点 $\vect{u},\vect{v}$ 成立。令 $m$ 足够大，得到
\[
  \abs{f(\vect{q}_m)-f(\vect{r}_m)}<\varepsilon,
\]
故两列极限相同。

若 $\vect{x}\in\QQ^n$，可取常值列 $\vect{q}_m=\vect{x}$，于是
\[
  F(\vect{x})=f(\vect{x}),
\]
故 $F$ 是扩张。

再证连续性，事实上 $F$ 仍一致连续。给定 $\varepsilon>0$，由原函数的一致连续性取 $\delta>0$，使对任意有理点 $\vect{u},\vect{v}$，
\[
  \abs{\vect{u}-\vect{v}}<\delta
  \Longrightarrow
  \abs{f(\vect{u})-f(\vect{v})}<\frac\varepsilon2.
\]
若
\[
  \abs{\vect{x}-\vect{y}}<\frac\delta3,
\]
分别选有理点列 $\vect{q}_m\to\vect{x}$、$\vect{r}_m\to\vect{y}$。当 $m$ 足够大时
\[
  \abs{\vect{q}_m-\vect{r}_m}<\delta,
\]
故
\[
  \abs{f(\vect{q}_m)-f(\vect{r}_m)}<\frac\varepsilon2.
\]
令 $m\to\infty$ 得
\[
  \abs{F(\vect{x})-F(\vect{y})}\le\frac\varepsilon2<\varepsilon.
\]
所以 $F$ 一致连续。

最后证惟一性。若 $G$ 是另一连续扩张，则对任意 $\vect{x}\in\RR^n$，取有理点列 $\vect{q}_m\to\vect{x}$，有
\[
  G(\vect{x})
  =\lim_{m\to\infty}G(\vect{q}_m)
  =\lim_{m\to\infty}f(\vect{q}_m)
  =F(\vect{x}).
\]
故 $G=F$，扩张惟一。
\end{xsolution}
